\section{Related Work}\label{related-work}

\subsection{System-Level Fault Diagnosis}
System-level diagnosis was introduced by Preparata, Metze, and Chien through the PMC model~\cite{pmc67}. Since then, many studies have investigated diagnosability, conditional diagnosability, and diagnosis algorithms for specific interconnection networks~\cite{wang1994diagnosability,zhu2016diagnosis,chang2020conditional,tang2024connectivity}. These works provide rigorous theoretical guarantees, but they often rely on deterministic testing assumptions and topology-specific structural properties. In contrast, this paper studies diagnosis under probabilistic testing behavior and aims to provide scalable posterior estimates rather than only deterministic fault sets.

\subsection{Probabilistic Diagnosis}
The probabilistic PMC model extends classical system-level diagnosis by modeling test outcomes as random variables~\cite{blough1989efficient,sun2023probabilistic}. This formulation is more realistic for large systems with noisy observations, but exact maximum-likelihood diagnosis is computationally expensive because the number of possible fault sets grows exponentially. Existing probabilistic diagnosis methods often require known model parameters or rely on restrictive assumptions about fault distributions. \tool differs by learning a graph-based approximation to node-wise posterior diagnosis and by explicitly considering calibration and parameter shifts.

\subsection{Graph Neural Networks for Diagnosis}
Graph neural networks learn representations by propagating information along graph edges~\cite{scarselli2008graph}. Representative architectures include GCN~\cite{kipf2017semi}, GraphSAGE~\cite{hamilton2017inductive}, and GAT~\cite{velivckovic2017graph}. These models are well suited to interconnection networks because diagnostic evidence is naturally local and relational. However, directly applying a GNN as a node classifier does not fully exploit the probabilistic structure of PMC diagnosis. \tool extends graph-based diagnosis by incorporating raw multi-round edge syndromes, learning edge reliability scores, and producing calibrated posterior probabilities instead of only hard labels.

\subsection{Uncertainty and Calibration}
Neural networks can be accurate but poorly calibrated, meaning that their predicted probabilities do not match empirical correctness. Calibration methods such as temperature scaling~\cite{guo2017calibration} and uncertainty estimation techniques such as Monte Carlo dropout~\cite{gal2016dropout} have been studied in machine learning. For fault diagnosis, calibration is particularly important because maintenance actions depend not only on which nodes are predicted faulty, but also on how reliable those predictions are. This paper brings calibration and uncertainty-aware ranking into probabilistic system-level diagnosis.